\subsection*{The geometry and its precise representation domain}
Here is the foundational calculation, including the choices needed to interpret it. Let $V=\mathbb R^4$ have a volume density. Real projective twistor space and its actual dual are $\mathbb P(V)$ and $\mathbb P(V^*)$. Compactified split conformal spacetime is $\operatorname{Gr}(2,V)$; a two-plane $E$ corresponds to the twistor line $\mathbb P(E)$ and the dual line $\mathbb P(E^\circ)$. Incidence is the natural pairing $W(Z)=0$. The conformal group is locally $SL(4,\mathbb R)\simeq\operatorname{Spin}_0(3,3)$, associated with spacetime signature $(2,2)$. This does not impose a metric of signature $(2,2)$ on $V$ itself.

For smooth even angular data, define homogeneous sections by
\[
 f(rZ)=|r|^{-2-w}f(Z),\qquad r\ne0.
\]
These are smooth density characters, not holomorphic line bundles of fractional degree. Initially, $-2<\Rea w<2$ gives the elementary locally integrable distribution interpretation on both sides. With Fourier kernel $e^{-2\pi iW(Z)}$ and compatible dual densities, distributional Fourier transformation gives
\[
 \mathcal F_w:\mathcal H_w(V)\longrightarrow\mathcal H_{-w}(V^*),\quad
 \mathcal F_w\pi_w(g)=\pi^*_{-w}(g)\mathcal F_w,
\]
where $\pi_w(g)f(Z)=f(g^{-1}Z)$ and $\pi^*_{-w}(g)h(W)=h(g^TW)$. The covariance follows by the substitution $Z=gY$, and ordinary Fourier inversion gives the even identity at generic parameters. Meromorphic homogeneous distribution continuation handles other parameters, with exceptional weights treated explicitly.

Choose an auxiliary Euclidean pairing only to identify the two compact pictures. Then the normalized operator has the twisted covariance
\[
 A(w)\pi_w(g)=\pi_{-w}(g^{-T})A(w).
\]
It is a generalized principal-series intertwiner for the projective maximal parabolic [4], rather than an untwisted commuting operator or a pointwise map between projective points. The harmonic spaces are $K=SO(4)$ types. In particular, the constant harmonic is $K$-fixed and is not a $G$-invariant subrepresentation at generic $w$.

At $w=0$, the projective scalar X-ray transform has the correct weight. In a matrix chart $Z=(\mu,\lambda)$ and $E_x=\{(x\lambda,\lambda)\}$, it is
\[
 (\mathcal Pf)(x)=\int_{\mathbb{RP}^1}f(x\lambda,\lambda)|\lambda\wedge d\lambda|.
\]
The two powers of the density cancel the degree $-2$. Differentiation under the integral gives
\[
 (\partial_{x_{11}}\partial_{x_{22}}-\partial_{x_{12}}\partial_{x_{21}})\mathcal Pf=0.
\]
since both integrands contain $\lambda_1\lambda_2\partial_{\mu_1}\partial_{\mu_2}f$. With compatible plane densities,
\[
 \mathcal P^*(\mathcal F_0f)(E^\circ)=\mathcal Pf(E).
\]
On a degree-$2m$ harmonic, complementary circle averages differ by $(-1)^m$, exactly the Fourier sign. Smooth harmonic expansion proves the identity. This classical split Penrose/X-ray statement [5] is asserted at the scalar weight $w=0$; the same unweighted map does not extend to all $w$ as one massless scalar representation.

\subsection*{Local factors and what the arithmetic completion chooses}
Put $\alpha=2+w$. Gaussian Mellin integration gives
\[
 \mathcal F(|Z|^{-\alpha})(W)=\gamma_\infty(\alpha)|W|^{\alpha-4},\qquad
 \gamma_\infty(\alpha)=\pi^{\alpha-2}\frac{\Gamma((4-\alpha)/2)}{\Gamma(\alpha/2)}.
\]
For a degree-$\ell$ spherical harmonic, the scalar is
$i^{-\ell}\pi^{\alpha-2}\Gamma((4+\ell-\alpha)/2)/\Gamma((\ell+\alpha)/2)$.
Dividing by the constant coefficient yields the product in \eqref{eq:model}.

At each prime, separately choose $\mathbb Q_p^4$, the max norm, the lattice $\mathbb Z_p^4$, a conductor-$\mathbb Z_p$ additive character, and Haar volume $\operatorname{vol}(\mathbb Z_p^4)=1$. At a unit dual vector, the shells with norm $p^k$, $k\leq0$, have trivial character; the shell $k=1$ has character integral $-1$, and later shells contribute zero. Thus
\[
 \gamma_p(\alpha)=(1-p^{-4})\sum_{k\leq0}p^{k(4-\alpha)}-p^{-\alpha}
 =\frac{1-p^{-\alpha}}{1-p^{\alpha-4}}.
\]
The unramified spherical product converges for $|\Rea w|<1$ and gives
\[
 C_4(w)=\gamma_\infty(2+w)\prod_p\gamma_p(2+w)
 =\frac{\Lambda(2-w)}{\Lambda(2+w)},\quad
 \Lambda(s)=\pi^{-s/2}\Gamma(s/2)\zeta(s).
\]
Selecting $B_4(w)=(w-1)(w-2)/((w+1)(w+2))$ to complete $\Lambda$ to $\xi$, and then using the global functional equation, gives $B_4C_4=R_{3/2}$. These arithmetic choices are extra data; covariance does not uniquely force them. No bounded adelic operator on all ramified types is claimed. Outside the initial strip the product is a meromorphically continued scalar identity, not a convergent cascade of local bounded maps.

Indeed, $\gamma_p(2+w)>1$ for every prime and every real $0<w<1$. The local finite-place spherical maps are expansive in this normalization. The completed wide bound uses the separate classical shifted-$\xi$ inequality [2], which is unconditional for every width $a\geq1/2$. Only $a=3/2$ has been obtained from the present four-dimensional local-factor construction. In the analogous dimension-$n$ calculation the completed width is $(n-1)/2$; analytic continuation in $n$ alone supplies no noninteger-dimensional twistor or $p$-adic space.

\subsection*{Bounded operators and distribution continuation are different claims}
For generic $w$, gamma-ratio asymptotics give
$b_m(w)=(-1)^m C(w)m^{-w}(1+O(m^{-1}))$ [7]. Hence $A(w)$ is compact and operator-norm holomorphic in $\CP$, with norm one. At positive even $w$, the product formula gives finite rank. Its unitary values on $i\mathbb R$ are strong-operator boundary limits, not norm limits: for $0<x<2$, $\|A(x)-A(0)\|=1$. At generic $\Rea w<0$, the multipliers grow, so $A(w)$ is not a bounded operator on the fixed $\HH$. Distributional Fourier inversion must not be read as bounded invertibility of this compact operator. The value $A(2)=P_0$ is a normalized continued limit at an exceptional raw homogeneous weight.

The chosen Euclidean compact-picture norm is $K$-invariant; the whole group action is unitary at imaginary $w$, not at arbitrary complex $w$. Thus the passive transfer interpretation refers to that specified norm, rather than a uniquely selected physical energy.
